\chapter{超対称性が重要でありうる理由}

\lectureoverview{真空揺らぎはこの章の諸主題を結びつける。境界条件はカシミール力を生み、ループ効果はヒッグス質量を脅かし、新たな仮想粒子は走るゲージ結合を変化させ、真空エネルギーは宇宙膨張を支配する。これらの問題は、ヒッグススケールを安定化させ、暗黒物質候補を供給し、ゲージ結合の統一を改善する方法としての低エネルギー超対称性を動機づける。次に我々はグラスマン座標を用いて超場を構成し、ベレジン積分を用いて超対称作用を抽出する。}

\section{反射板の間の真空揺らぎ}

我々はまずカシミール効果に関する一つの問いを片づけることによって超対称性に戻る。二枚の平行な反射板を距離\(d\)だけ離して置く。実粒子が存在しないときでさえ、量子場には真空揺らぎがある。それらを仮想過程として、あるいは同等に、場の振動子モードの零点運動として思い描いてよい。板は境界条件を課すので、それらの間で許される揺らぎは、制限のない空間における揺らぎとは異なる。

重要な量はエネルギーの絶対的な零点ではなく、間隔が変化するときに生じる変化である。単位面積あたりの真空エネルギーを\(\mathcal E(d)=E(d)/A\)で表す。\(\mathcal E\)が\(d\)に依存するようになれば、それはポテンシャルエネルギーであり、単位面積あたりの力は
\begin{equation}
\frac{F}{A}=-\frac{d\mathcal E}{dd}.
\end{equation}
である。したがって、仮想過程の再配置が観測可能な力となる。理想導体板というよく知られた電磁気の例では、エネルギーと圧力は
\begin{align}
\mathcal E(d)&=-\frac{\pi^2}{720\,d^3},\\
\frac{F}{A}&=-\frac{\pi^2}{240\,d^4},
\end{align}
であり、単位は\(\hbar=c=1\)である。\editorialnote{これらの係数と、\(d^{-3}\)の単位面積あたりのエネルギーと\(d^{-4}\)の圧力との区別が、黒板で提起された次元の問題を解決する。それらは標準的な理想板の結果であり、ここで行った導出ではない。} 負号は引力を表す。

\begin{classroomqa}
\audiencequestion{4乗は力を表すのか、それともエネルギーを表すのか。\lecturetimestamp{00:04:39}}
\lecturerresponse{4次元時空における平行板では、単位面積あたりのエネルギーは\(d^{-3}\)で減少し、一方\(d\)に関して微分すると単位面積あたりの力は\(d^{-4}\)で減少する。\(\hbar\)と\(c\)を復元すれば、次元解析だけでこれらの冪はすでに確定する。}
\end{classroomqa}

板が仮想光子ループを遮ると言って計算を要約することがある。その描像は、あまり文字どおりに受け取らない限り有用である。正確な計算は、許される場のモードのスペクトルの変化、あるいは同等に、量子有効作用の変化である。どちらの言葉遣いでも同じ間隔依存のエネルギーに至る。

\begin{classroomqa}
\audiencequestion{カシミール力は引力ではなく斥力になりうるか。\lecturetimestamp{00:07:48}}
\lecturerresponse{なりうる。符号は普遍的ではなく、場、幾何、物質、境界条件に依存する。フェルミオンループはしばしばボソンループとは反対のループ符号をもたらし、それが黒板でなされた示唆を動機づけるが、フェルミオン統計それ自体は、すべてのカシミール力が斥力でなければならないという一般的定理ではない。\editorialnote{この限定は、意図的にためらいがちな教室での答えを、誤った普遍的規則を避けるために必要な一般的陈述に置き換える。}}
\end{classroomqa}

\subsection{真空エネルギーの証拠として何が数えられるか。}

カシミール力は測定されており、原子の観測量も真空ループに敏感である。しかし、どちらの事実もそれ自体では、絶対的な真空エネルギーが重力を及ぼすことを示すものではない。カシミール実験は境界によって引き起こされるエネルギー差を測定する。同じ力は、非有界な絶対零点和に操作的な意味を与えることなく記述できる。

決定的な検証は重力から来る。エネルギーは重力荷である。もし空の空間が一様なエネルギー密度を担うなら、それはアインシュタイン方程式に入り、宇宙膨張に影響を与えなければならない。したがって、宇宙の加速膨張は、真空的な成分に対する直接の重力的証拠である。

\begin{classroomqa}
\audiencequestion{なぜ一般向けの解説は、物理学者がそれが何であるかを知らないかのように、真空エネルギーを神秘的なものと呼ぶのか。\lecturetimestamp{00:08:19}}
\lecturerresponse{量子場の理論はすべての場に零点寄与を与えるので、真空エネルギーの可能性は神秘的ではない。神秘は定量的である。なぜ重力を及ぼす値は、既知の量子場と高エネルギーカットオフに基づく推定に比べて、これほど並外れて小さいのか。驚くべき事実は、それがほとんど存在しないことであり、単に存在することではない。}
\end{classroomqa}

長年にわたり、観測による上限はあまりに小さかったため、宇宙定数を厳密にゼロに置き、その相殺を強制する原理を探すことが心理的に誘惑的であった。しかし改良された宇宙観測は、代わりに小さな非ゼロの値を発見した。この発見によって真空エネルギーが概念的に異質になったわけではない。説明されない抑制の度合いが無視できなくなったのである。関連する比較は仮定する紫外スケールに依存するが、自然な高エネルギー推定はどれも巨大な相殺を要求する。

\section{多粒子状態に一つの超電荷が及ぼす効果}

超対称性を信じる理由を列挙する前に、可能なパラドックスを取り除いておくべきである。ボソンは一つのモードを任意の大きな数で占有しうるが、フェルミオン的モードの占有数はゼロか1である。もし対称性がボソンとフェルミオンを交換するなら、一つのモードにある千個のボソンを千個の禁じられたフェルミオンに変えるのだろうか。

いいえ。フェルミオン的生成子は一度に一つの量子を変える。振動子模型で、\(a\) がボソンを消滅させ、\(c^\dagger\) がそのフェルミオン的パートナーを生成するとしよう。一つの超電荷は次のように表されうる。
\begin{equation}
Q=c^\dagger a.
\end{equation}
\(n\) 個のボソンを持ちフェルミオンを持たない状態に作用させると、
\begin{equation}
Q\ket{n,0}
=\sqrt n\,\ket{n-1,1}.
\end{equation}
一つのボソンが取り除かれ、一つのフェルミオンが挿入された。同じ電荷で再び作用させるとゼロを与える。
\begin{equation}
Q^2=(c^\dagger)^2a^2=0.
\end{equation}
これはフェルミオン的生成演算子から受け継がれた冪零性である。

\begin{classroomqa}
\audiencequestion{多数のボソンが同じ状態を占めるとき、なぜ超対称性変換はその同じ状態に多数のフェルミオンを生成して排他原理に反しないのか。\lecturetimestamp{00:12:05}}
\lecturerresponse{生成子は集団全体を置き換えるのではない。状態をフェルミオン数の1単位だけ変える。同じフェルミオン的電荷は、その二乗が消えるため、通常の回転生成子のようには合成できない。相補的な電荷が交換を元に戻すことができ、混合反交換子は時間並進で閉じる。}
\end{classroomqa}

最後の文は本質的である。略記 \(Q^2=0\) は個々のカイラル成分、または今書いた単純な電荷に適用されるのであり、超電荷のあらゆる積が消えるという意味ではない。四次元超対称性では、
\begin{equation}
\{Q_\alpha,\bar Q_{\dot\beta}\}
=2\sigma^\mu_{\alpha\dot\beta}P_\mu,
\end{equation}
一方、中心電荷のない最小代数に対しては
\begin{equation}
\{Q_\alpha,Q_\beta\}=0,
\qquad
\{\bar Q_{\dot\alpha},\bar Q_{\dot\beta}\}=0
\end{equation}
である。\editorialnote{この成分形は、同じ \(Q\) の繰り返し作用が消え、超電荷とその共役が並進を生成するという教室での略記を正確にする。}

\section{低エネルギー超対称性を探す三つの理由}

\begin{classroomqa}
\audiencequestion{真空エネルギーを相殺すること以外に超対称性の動機はあるか。\lecturetimestamp{00:15:06}}
\lecturerresponse{密接に関連する三つの動機がある。ヒッグス自己エネルギーの制御、安定な暗黒物質粒子の可能性、そして高エネルギーにおけるゲージ結合定数のはるかに鋭い一致である。いずれも超対称性が自然界で実現していることを証明するものではないが、合わせて考えると、検証すべき真剣な理論にしている。}
\end{classroomqa}

\emph{低エネルギー超対称性} という語句は相対的である。超対称パートナーが日常的な尺度で軽いという意味ではなく、その質量が電弱スケールや加速器スケールをはるかに超えていないという意味である。超対称パートナーが任意に高いところへ押しやられれば、三つの論拠はすべて弱まる。

\subsection{ヒッグススケールの保護}

スカラー質量は高エネルギー物理に異常に敏感である。模式的に、ボソン的ループは次の形の補正を生む。
\begin{equation}
\delta m_H^2\big|_{\mathrm B}
\sim +\frac{\lambda}{16\pi^2}\Lambda^2,
\end{equation}
ここで \(\Lambda\) は有効理論が用いられる上限のスケールを表す。フェルミオン的ループは反対の符号を持つ。
\begin{equation}
\delta m_H^2\big|_{\mathrm F}
\sim -\frac{|y|^2}{8\pi^2}\Lambda^2.
\end{equation}
係数は相互作用と場の多重度に依存するが、超対称性はそれらを関係づけ、主要な紫外感度がパートナー間で相殺されるようにする。\editorialnote{これらの模式的なループ公式は、定性的な自己エネルギー論の背後にある数学を明示するものである。標準模型のヒッグス質量の完全な計算を意図したものではない。}

これは重要である。なぜならヒッグス期待値が電弱スケールを定めるからである。ゲージボソンとフェルミオンの質量は次のような関係に従う。
\begin{equation}
m_W=\frac{gv}{2},
\qquad
m_f=\frac{y_fv}{\sqrt2}.
\end{equation}
もしヒッグスセクターが自然に巨大なスケールへ押しやられるなら、裸のパラメータが輻射補正に対して極めて精密に調整されない限り、よく知られた低エネルギー粒子スペクトルは失われるであろう。

厳密な超対称性は相殺を厳密にするであろうが、厳密な超対称性は観測されたボソンとフェルミオンの質量が等しくないことによってすでに排除されている。ソフトに破れた超対称性では、二次感度は超対称パートナーの分裂への依存に置き換えられ、模式的に
\begin{equation}
\delta m_H^2
\sim
\frac{g^2}{16\pi^2}
\left(m_{\mathrm{soft}}^2-m_{\mathrm{ordinary}}^2\right)
\ln\frac{\Lambda^2}{m_{\mathrm{soft}}^2}.
\end{equation}
となる。相殺が有用であり続けるのは、分裂が大きすぎない場合に限られる。

\subsection{暗黒物質カスケードの安定な終端点}

スペクトルを通常の粒子とその超対称パートナーに分ける。フェルミオンにはボソン的パートナーがあり、ボソンにはフェルミオン的パートナーがある。相互作用がこれら二つのクラスを区別するパリティを保存するならば、超対称パートナーの崩壊は別の超対称パートナーを残さなければならない：
\begin{equation}
\widetilde X
\longrightarrow X+\widetilde Y
\longrightarrow X+Y+\widetilde Z
\longrightarrow\cdots.
\end{equation}
このカスケードは最後の超対称パートナーを捨てることはできない。したがって、それは最も軽い超対称粒子、すなわちLSPで終わる。

よく用いられる対称性は\(R\)パリティであり、
\begin{equation}
R_p=(-1)^{3(B-L)+2s}.
\end{equation}
標準模型の粒子は\(R_p=+1\)を持ち、その超対称パートナーは\(R_p=-1\)を持つ。\(R_p\)が厳密ならば、すべての相互作用は偶数個の超対称パートナーを含み、LSPは安定である。\editorialnote{超対称性だけでは安定なLSPは保証されない。安定性には保存される\(R\)パリティ、または同じ帰結をもつ別の対称性が必要であり、それがなければLSPは崩壊しうる。}

暗黒物質は少なくとも宇宙論的な時間スケールにわたって生き残らなければならない。また、可視の銀河円盤を繰り返し通過しても実質的なエネルギーを失わない程度に弱くしか相互作用してはならない。電荷をもつ、あるいは強い相互作用をする集団は散乱し、エネルギーを散逸させ、通常の物質と同じ薄い構造へと崩壊する傾向がある。それに対して観測される広がったハローは、中性で弱く相互作用する成分を示唆している。

これはLSPを一意に同定するものではない。模型と対称性の破れ方に応じて、候補は電弱ゲージーノとヒグシーノの中性混合状態かもしれないし、グラビティーノかもしれない。電荷をもつ候補やカラーをもつ候補は強く制限される。超対称性が与えるものは保証された答えではなく、中性で長寿命の終端点を生成する自然な機構である。

\section{走る結合定数と大統一の兆候}

第三の動機は電荷の遮蔽から出発する。測定される結合定数がエネルギー・スケールに依存するのは、仮想粒子が相互作用を着飾るからである。電磁気学では、光子伝播関数に挿入された電子・陽電子ループが長距離で電荷を遮蔽する。より短い距離、すなわちより高い運動量移行では、その分極雲のうちプローブの間にある部分が少なくなり、有効電荷はより大きくなる。

次のように書く：
\begin{equation}
\alpha_i=\frac{g_i^2}{4\pi}.
\end{equation}
1ループでは、逆結合定数はエネルギーの対数に対して線形に走る：
\begin{equation}
\frac{d\alpha_i^{-1}}{d\ln\mu}
=-\frac{b_i}{2\pi}.
\end{equation}
したがって
\begin{equation}
\alpha_i^{-1}(\mu)
=\alpha_i^{-1}(\mu_0)
-\frac{b_i}{2\pi}\ln\frac{\mu}{\mu_0}.
\end{equation}
\editorialnote{この1ループの式は、対数エネルギーを変数として描かれたほぼ直線の正確な版である。閾値と高次ループの補正がその直線をわずかに曲げる。}

電弱スケールより上では、三つの基本的なゲージ結合定数は
\begin{equation}
U(1)_Y,\qquad SU(2)_L,\qquad SU(3)_c.
\end{equation}
のものである。低エネルギーで直観的に用いられる電磁結合定数は、電弱対称性の破れの後では最初の二つの組み合わせである。QCDでは、非可換ゲージ場が反遮蔽を生み出す：強い結合定数は短距離で減少する。これが漸近的自由性という現象である。その結果、\(\alpha_3^{-1}\)は\(\ln\mu\)とともに増大するが、他の逆結合定数は異なる傾きをもつ。

標準模型の曲線は互いに示唆的に近づくが、一点で交わらない。超対称パートナーは仮想ループを変更し、したがって係数\(b_i\)を変える。従来の大統一の規格化では、
\begin{align}
(b_1,b_2,b_3)_{\mathrm{SM}}
&=\left(\frac{41}{10},-\frac{19}{6},-7\right),\\
(b_1,b_2,b_3)_{\mathrm{MSSM}}
&=\left(\frac{33}{5},1,-3\right).
\end{align}
超対称パートナーの閾値がTeVスケール付近にあると、外挿された超対称曲線は
\begin{equation}
M_{\mathrm{GUT}}\sim 10^{16}\ \mathrm{GeV}.
\end{equation}
付近ではるかに正確に交わる。
どの実験もそのスケールに直接到達しない。低エネルギーの結合定数が測定され、場の内容が走りを決定し、曲線は何十年にもわたって外挿される。二本の線はどこかで必ず交わる。三本の独立な線が一点付近で交わることこそが非自明な観測である。

\begin{classroomqa}
\audiencequestion{改善された外挿は厳密な超対称性と等しいパートナー質量を仮定しているのか？\lecturetimestamp{00:38:09}}
\lecturerresponse{いいえ。超対称性は破れているので、質量は分裂している。重要なのは、超対称パートナーのループがいつ走りに寄与し始めるかである。パートナーがあまり重くなければ、傾きは十分早く変化し、三つの結合定数は収束する。それらが極端に重ければ、超対称の閾値は遅すぎる時期に現れ、改善された一致は失われる。}
\end{classroomqa}

閾値補正と未知の高エネルギー・スペクトルのために、三本の線が数学的に厳密に交わると文字通り主張することはできない。それでも結果は際立っている：ヒッグスを助け、熱的暗黒物質候補を与えるのと同じおおよその超対称パートナー・スケールが、大統一も改善するのである。

\subsection{重力と共通の高エネルギー・スケール}

重力はその電荷がエネルギーそのものであるため、異なる振る舞いをする。重心系エネルギー\(E\)における有用な無次元の尺度は
\begin{equation}
\alpha_G(E)\sim GE^2
=\left(\frac{E}{M_{\mathrm{Pl}}}\right)^2.
\end{equation}
である。通常のエネルギーではこれは極めて小さいが、\(E\)とともに二次関数的に増大する。したがってその逆数は、対数的に走るゲージ結合の逆数よりもはるかに速く減少し、プランクスケール付近で1の程度になる。

この比較は、すべての相互作用が近接した高いスケールによって支配されている可能性を示唆するが、ゲージの統一ほどきれいではない。ゲージ結合は繰り込み可能な場の理論における無次元パラメータである。ニュートンの定数は次元をもつため、重力を同じプロット上に移すには規約が必要である。

\begin{classroomqa}
\audiencequestion{超対称パートナーは同じ3つのゲージ相互作用に参加するのか、それとも新しい力を定義するのか。\lecturetimestamp{00:42:35}}
\lecturerresponse{それらはパートナーと同じゲージ量子数をもつ。スレプトンは電子の電荷をもち、スクォークは対応するクォークの色電荷と電弱電荷をもつが、スピンと質量は異なる。したがってそれらのループは、3つの新しい相互作用を導入するのではなく、同じゲージ相互作用を修正する。}
\end{classroomqa}

\begin{classroomqa}
\audiencequestion{LHCで超対称パートナーが発見されれば、超対称性の決定的な証拠になるだろうか。\lecturetimestamp{00:43:27}}
\lecturerresponse{一つの未知の粒子だけでは、その名前だけでは十分ではないが、スピン、電荷、質量、結合の間の超対称的な関係に従う超対称パートナーのスペクトルは非常に特徴的であろう。超対称性は多くの相関した制限を課すので、詳細な相互作用のパターンが決定的な検証を与えるだろう。}
\end{classroomqa}

\subsection{なぜスケールを犠牲なしに遠ざけることができないのか}

超対称パートナーが実験の到達範囲のすぐ先にしかないなら、3つの動機はすべて依然として働きうる。それらの質量が桁違いに引き上げられると、それぞれの議論は悪化する：
\begin{enumerate}
\item ヒッグス質量の相殺は、分裂に対する大きな残留感度を残す。
\item 従来の弱いスケールの熱的残存粒子は、もはや期待されるスケールに現れない。
\item ゲージ結合の走りの閾値が上方に移動し、3つの結合はより不正確にしか交わらない。
\end{enumerate}
超対称性は依然としてはるかに高いスケールに存在しうるが、もはやこれらの特定の低エネルギー問題を同じ経済的な方法で解決することはないだろう。このアイデアの長寿命性は、自然がそれを用いなければならない証拠ではない。それは、この提案が決定的な検証に値するほど首尾一貫している証拠である。

\section{宇宙論が暗黒物質について教えうることと教ええないこと}

\begin{classroomqa}
\audiencequestion{標準宇宙論は、元素合成が水素とヘリウムを制限するのと同じように、暗黒物質の量と各暗黒物質粒子の質量を決定できるだろうか。\lecturetimestamp{00:46:32}}
\lecturerresponse{宇宙論は暗黒物質の全質量密度を非常によく測定するが、密度だけでは質量と数を分離できない。同じ密度は、少数の重い粒子によって担われることも、膨大な数の非常に軽い粒子によって担われることもある。微視的な粒子を同定するには、追加の理論または非重力観測が必要である。}
\end{classroomqa}

二つの広い可能性がこの曖昧さを例示する。電弱スケールまたはTeVスケール付近の弱い相互作用をする粒子は、高温の熱浴から凍結脱出しうる。もう一方の極端では、アクシオンのような極めて軽いボソンがコヒーレントな場の配置または凝縮体を占めうる。中間の可能性は、二つの最もよく発達した機構がほぼ両極端に位置するというだけの理由で禁じられるわけではない。

ハロの構造自体が重要な制限を与える。暗黒物質は通常の物質とほとんど摩擦なしに交差するため、空間的に広がったままであり、輝く円盤に崩壊しない。回転曲線やその他の重力現象は質量が存在することを教えるが、それ自体では粒子を数え上げることはない。

\begin{classroomqa}
\audiencequestion{初期宇宙は当初、超対称パートナーのみを含み、それが後に我々が見る粒子へ崩壊した可能性はあるだろうか。\lecturetimestamp{00:49:11}}
\lecturerresponse{すべての関連する質量よりはるかに高い温度では、急速な反応が種を任意の一方的な集団ではなく熱平衡へと駆動する。通常の衝突は超対称パートナーが不足しているときにそれらを生成し、超対称パートナーの衝突と崩壊は通常の粒子を補充する。正確な存在量は相互作用と自由度に依存するが、高温プラズマは絶えず釣り合いへと向かう。}
\end{classroomqa}

宇宙が冷えるにつれて、重い種はボルツマン抑制を受け、対消滅または崩壊する。安定なLSPは、その相互作用率が宇宙膨張より遅くなると、もはや平衡状態を追跡しなくなる：
\begin{equation}
\Gamma_{\mathrm{ann}}
\sim n\langle\sigma v\rangle
\lesssim H.
\end{equation}
残った粒子は残存存在量を形成する。\editorialnote{この凍結脱出条件は、熱平衡の議論に対する標準的な定量的完成である。他のLSP候補、特に非常に弱く結合したグラビティーノは、非熱的な履歴を持ちうる。}

\begin{classroomqa}
\audiencequestion{LSPはどの通常粒子のパートナーであり、それは非相互作用でなければならないのか？ \lecturetimestamp{00:51:05}}
\lecturerresponse{モデルに依存しない答えは存在しない。実行可能な暗黒物質LSPは電気的に中性で無色であると期待されるが、完全に非相互作用である必要はない。それは光子、弱ゲージボソン、およびヒッグス場に関連する混合状態であってもよいし、グラビティーノであってもよい。その正確な正体は、超対称性の破れのスペクトルから従う。}
\end{classroomqa}

\section{真空図、自己エネルギー図、および膨張}

\begin{classroomqa}
\audiencequestion{ヒッグス発散と真空エネルギー発散の違いは何か？ \lecturetimestamp{00:52:09}}
\lecturerresponse{真空図には外部粒子線がなく、空の空間の有効作用を変化させ、真空エネルギーに寄与する。ヒッグス自己エネルギー図には2本の外部ヒッグス脚があり、ヒッグス伝播関数を変化させ、したがってその質量を変化させる。どちらも仮想揺らぎを含むが、それらは異なる問いに答える。}
\end{classroomqa}

実粒子を、真空に挿入された不純物として思い描いてもよい。この不純物に付着する真空揺らぎは、その伝播と質量をずらす。カシミール板も真空揺らぎを変更するが、それは巨視的な境界条件を通じてである。共通の要素は、挿入されたものに対する量子揺らぎの応答であり、補正される観測量は各場合で異なる。

\vspace{-2pt}
\subsection{物質、放射、および真空は異なる仕方で薄まる}

\(a(t)\)を宇宙スケール因子、
\begin{equation}
H=\frac{\dot a}{a}
\end{equation}
をハッブルパラメータとする。一様な宇宙における局所的な応力エネルギーの保存は
\begin{equation}
\dot\rho+3H(\rho+p)=0.
\end{equation}
を与える。状態方程式\(p=w\rho\)に対して、これは
\begin{equation}
\rho(a)\propto a^{-3(1+w)}.
\end{equation}
と積分される。
\editorialnote{連続の式は、粒子数え上げおよび波長伸長の説明の根底にある簡潔な導出を与える。}

冷たい非相対論的物質では\(p\approx0\)であるから、
\begin{equation}
\rho_m\propto a^{-3}.
\end{equation}
固定された共動体積は\(a^3\)として成長する一方、その粒子が運ぶ静止エネルギーはほぼ固定されたままである。

放射では\(p=\rho/3\)であり、
\begin{equation}
\rho_r\propto a^{-4}.
\end{equation}
3つの冪は数密度の希薄化に由来する。4つ目は赤方偏移に由来する：
\begin{equation}
\lambda\propto a,
\qquad
E_\gamma=\frac{2\pi\hbar c}{\lambda}\propto a^{-1}.
\end{equation}

\begin{classroomqa}
\audiencequestion{光子が波動関数を持つなら、なぜ希薄化則を導出する際にそれを波ではなく粒子として語ってよいのか？ \lecturetimestamp{00:58:21}}
\lecturerresponse{この区別はスケーリングに影響しない。粒子描像では、共動体積内の数は固定されており、各量子は\(a^{-1}\)として赤方偏移する。古典波描像では、波長が伸び、場のエネルギーは同じ量だけ減少する。膨張を見るのが最も容易になる言語をどちらでも使うことができる。}
\end{classroomqa}

真空エネルギーは\(p_\Lambda=-\rho_\Lambda\)であるから、
\begin{equation}
\rho_\Lambda=\text{constant}.
\end{equation}
空間が膨張しても薄まらない。実際、共動領域は
\begin{equation}
E_\Lambda=\rho_\Lambda a^3V_0,
\end{equation}
を含み、これはその物理的体積とともに成長する。この事実が、我々がまもなく扱うエネルギー保存の問題を生み出す。

異なる冪が宇宙の歴史を組織する。放射は最も速く減少し、物質はよりゆっくり減少し、真空エネルギーはまったく減少しない。したがって、微小な真空項は初期宇宙では無視でき、後期には支配的になりうる。ここで議論した概略値を用いると、現在のエネルギー収支は約\(70\%\)の真空エネルギーと\(30\%\)の物質であり、放射ははるかに小さい。

\section{幾何学はエネルギーに等しい}

アインシュタイン方程式の時間-時間成分は、前述のスケーリング則を力学に変える。黒板はスローガンと1つの幾何学的項から始まる。

\begin{figure}[H]
\centering
\includegraphics[width=0.68\textwidth]{lecture_06_figure_03.png}
\lectureframe{01:06:37.088}
\caption{組織化するスローガン\emph{幾何学はエネルギーに等しい}が、その下に膨張項\(H^2=(\dot a/a)^2\)が書かれている。}
\end{figure}

一様で等方な宇宙に対して、第一フリードマン方程式は
\begin{equation}
H^2+\frac{k}{a^2}
=\frac{8\pi G}{3}\rho,
\end{equation}
単位\(c=1\)で表され、ここで\(k=+1,0,-1\)は閉じた、平坦な、または開いた空間幾何学を表し、
\begin{equation}
\rho
=\rho_{\mathrm{ordinary}}
+\rho_{\mathrm{dark}}
+\rho_{\mathrm{radiation}}
+\rho_{\mathrm{vacuum}}.
\end{equation}
である。\(1/a^2\)の隣に見られる符号は、曲率項が左辺に保たれるか右辺に移されるか、および\(k\)の値に依存する。

\begin{figure}[H]
\centering
\includegraphics[width=0.96\textwidth]{lecture_06_figure_04.png}
\lectureframe{01:08:46.794}
\caption{完全な板書の収支：幾何学的側面には膨張と空間曲率が、源の側面には物質、放射、真空エネルギーが入る。物質はさらに通常成分と暗黒成分に分けられる。}
\end{figure}

\subsection{宇宙論的ハミルトニアン拘束条件}

\begin{classroomqa}
\audiencequestion{宇宙が膨張するにつれて共動領域内の真空エネルギーが増大するなら、正味のエネルギー密度は一定なのか、それとも全エネルギーが保存されるのか。\lecturetimestamp{01:01:51}}
\lecturerresponse{一般相対性理論は、任意の膨張時空に対して普遍的な大域的エネルギーを与えるものではない。しかし、一様な宇宙論的縮約においては、時間-時間アインシュタイン方程式はハミルトニアン拘束条件である。幾何学と物質を一方の側に置くと、拘束されたハミルトニアンはゼロになる。変化する物質の項は、箱の中の通常の固定されたエネルギーによってではなく、重力自由度によって釣り合わされる。}
\end{classroomqa}

概略的に、正準的な主張は次式で表される。
\begin{equation}
\mathcal H_{\mathrm{gravity}}
+\mathcal H_{\mathrm{matter}}
=0.
\end{equation}
これが、全エネルギーがゼロであるという言明の制御された意味である。これは時間再パラメータ化不変性によって課される運動方程式であり、すべての局所エネルギー密度がゼロになるという主張ではない。共動体積において、通常の物質エネルギーはほぼ一定であり、放射エネルギーは減少し、真空エネルギーは増大する。一方、重力の寄与と膨張率が調整され、アインシュタインの拘束条件が満たされ続ける。

局所保存は厳密に成り立つ：
\begin{equation}
\nabla_\mu T^{\mu\nu}=0.
\end{equation}
一様流体に対して、これはまさに連続の方程式
\(\dot\rho+3H(\rho+p)=0\)である。一般に破れるのは、膨張宇宙において物質エネルギーの座標に依存しない大域的総和も保存されなければならないという平坦時空の直観である。\editorialnote{この区別により、エネルギーは常にゼロであるという教室での言明が正確になる：局所共変保存と宇宙論的ハミルトニアン拘束条件は厳密であるが、一意な大域的物質＋重力エネルギーは一般には定義されない。}

\begin{classroomqa}
\audiencequestion{ゼロ拘束条件は、宇宙が開いているか、閉じているか、平坦かを決める条件なのか。\lecturetimestamp{01:04:42}}
\lecturerresponse{いいえ。ハミルトニアン拘束条件はすべての場合に適用される。開いた宇宙、平坦な宇宙、閉じた宇宙は曲率ラベル\(k\)によって、したがって\(k/a^2\)の項によって異なる。曲率は収支への一つの寄与を変えるが、拘束条件をオンにしたりオフにしたりするものではない。}
\end{classroomqa}

\section{グラスマン座標と超場}

\begin{classroomqa}
\audiencequestion{グラスマン数は実際に理論のラグランジアンに現れるのか。\lecturetimestamp{01:09:37}}
\lecturerresponse{それらはまず、フェルミオン場がグラスマン奇であるために現れる。次に、超空間の座標として再び現れる。超場はそれらの奇座標に依存し、超対称な作用はそれらについて微分し積分することによって構成される。}
\end{classroomqa}

通常の場は時空の関数である：
\begin{equation}
\phi=\phi(x^\mu).
\end{equation}
四次元の最小超対称性は、\(x^\mu\)を二成分奇スピノル\(\theta^\alpha\)とその共役\(\bar\theta^{\dot\alpha}\)によって拡張する：
\begin{equation}
z^M=(x^\mu,\theta^\alpha,\bar\theta^{\dot\alpha}).
\end{equation}
二つの複素奇成分、すなわち四つの実奇方向が存在する。超場はこの超空間上の関数である：
\begin{equation}
\Phi=\Phi(x,\theta,\bar\theta).
\end{equation}

これらは計量の意味での微小な空間次元ではない。グラスマン座標には正の数値的な大きさがない。それを小さいと呼ぶのは、冪零性のための記憶術である：
\begin{equation}
(\theta^\alpha)^2=0,
\qquad
\theta^\alpha\theta^\beta=-\theta^\beta\theta^\alpha.
\end{equation}
\editorialnote{有限個のグラスマン座標におけるあらゆる展開が終わる理由は、数値的な小ささではなく冪零性である。}

四つの独立な奇座標があるとき、完全に一般的な超場は最大\(2^4=16\)個のグラスマン単項式を持つ。概略的に、
\begin{equation}
\begin{aligned}
\Phi(x,\Theta)={}&\phi(x)+\Theta_A\psi_A(x)\\
&+\frac{1}{2}\Theta_A\Theta_B Z_{AB}(x)+\cdots\\
&+\Theta_1\Theta_2\Theta_3\Theta_4D(x),
\end{aligned}
\end{equation}
ここで\(\Theta_A\)は四つの実奇座標をまとめて表す。超場全体を偶に選べるので、偶数乗の係数はボソン的であり、奇数乗の係数はフェルミオン的である。したがって一つの超場は、通常の成分場の有限個の多重項をまとめたものである。

このまとめ方は単なる記法ではない。超対称性は、時空の並進と混ざり合った奇方向の並進として幾何学的に作用する。超空間の規則に従って超場から構成されたラグランジアンは、テンソルから構成されたスカラーラグランジアンが回転不変性やローレンツ不変性を受け継ぐのと同じように、その対称性を受け継ぐ。

作用は今やグラスマン測度を含む。典型的な全超空間の項は次の形をとる：
\begin{equation}
S_D=\int d^4x\,d^2\theta\,d^2\bar\theta\,
\mathcal K(\Phi,\Phi^\dagger).
\end{equation}
一方、カイラル項は次を用いる：
\begin{equation}
S_F=\int d^4x\,d^2\theta\,W(\Phi)+\mathrm{h.c.}
\end{equation}
\editorialnote{全超空間測度とカイラル測度の区別は、概略的な\(d^4x\,d^4\theta\)作用の標準的な精密化である。どちらの測度が適切かは具体的な例が決定する。}

したがって、通常の微積分では答えられない問いに到達した。反可換座標にわたる積分とは何を意味するのか。

\section{Berezin積分}

\subsection{通常のモデル：全微分}

大きな正および負の\(x\)で消える通常の関数\(F(x)\)から始めよう。微積分学の基本定理により
\begin{equation}
\int_{-\infty}^{\infty}dx\,\frac{dF}{dx}
=F(\infty)-F(-\infty)
=0.
\end{equation}

\begin{figure}[H]
\centering
\includegraphics[width=0.88\textwidth]{lecture_06_figure_05.png}
\lectureframe{01:26:31.594}
\caption{Berezin積分の通常の原型：全微分の積分は端点の値に帰着し、関数が減衰すれば消える。黒板の下の行は一部隠れているため、きれいな式を本文に示す。}
\end{figure}

Grassmann積分が線形であり、全微分がゼロに積分されるという有用な性質を共有することを望む。1つの奇座標に対して、次を要求する。
\begin{align}
\int d\theta\,[af(\theta)+bg(\theta)]
&=
a\int d\theta\,f(\theta)
+b\int d\theta\,g(\theta),\\
\int d\theta\,\partial_\theta f(\theta)&=0.
\end{align}
第2の条件は、測度の並進不変性の代数的類似である。奇数の数直線には端点がなく、微小間隔にわたる和もない。

\begin{classroomqa}
\audiencequestion{不定Grassmann積分が存在しないなら、積分は微分の逆であるという考えを放棄するのか。\lecturetimestamp{01:21:19}}
\lecturerresponse{はい。Berezin積分は定値線形汎関数であり、原始関数を求める操作ではない。1つのGrassmann変数に対しては、実際、微分と同じ代数的操作で表される。}
\end{classroomqa}

\begin{classroomqa}
\audiencequestion{全時空にわたる\(d^4x\)の積分もまた定積分なのか。\lecturetimestamp{01:21:34}}
\lecturerresponse{はい。作用を時空にわたって積分するとき、境界条件は了解されている。違いは、通常の時空積分は和の極限を通じて理解できるのに対し、Grassmann積分は有限基底\(1,\theta\)への作用によって代数的に定義されることである。}
\end{classroomqa}

\subsection{1変数}

1つの奇変数の任意の関数は、その係数を固定された右側の順序で書ける。
\begin{equation}
f(\theta)=a+\theta b.
\end{equation}
\(\partial_\theta\theta=1\)であるから、全微分の規則は次を意味する。
\begin{equation}
\int d\theta\,1=0.
\end{equation}
残りの規格化は次に選ばれる。
\begin{equation}
\int d\theta\,\theta=1.
\end{equation}
したがって、
\begin{equation}
\int d\theta\,f(\theta)=b.
\end{equation}
積分は利用可能な最高次の係数を選び出す。ここで用いる左微分の規約では、
\begin{equation}
\partial_\theta f=b
=\int d\theta\,f.
\end{equation}

\begin{classroomqa}
\audiencequestion{\(\partial_\theta F=\theta\)を満たす関数\(F(\theta)\)は存在するか。\lecturetimestamp{01:30:05}}
\lecturerresponse{いいえ。そのような関数には\(\theta^2\)に比例する項が必要だが、\(\theta^2=0\)である。これは、不定原始関数のGrassmann類似が存在しない理由を別の形で見る方法である。}
\end{classroomqa}

\subsection{2変数と測度の順序}

複素対\(\theta,\bar\theta\)に対して、一般の偶関数を順序付けられた形で書く。
\begin{equation}
f(\theta,\bar\theta)
=a+\bar\theta\psi+\bar\psi\theta+b\,\bar\theta\theta.
\end{equation}
\begin{equation}
\theta^2=\bar\theta^2=0,
\qquad
\bar\theta\theta=-\theta\bar\theta.
\end{equation}であるから、より高次の項は存在しない。

ここで規約を固定する。
\begin{equation}
\int d\theta\,d\bar\theta\,
\bar\theta\theta=1.
\end{equation}
反復積分では、隣接する\(d\bar\theta\)の操作が最初に作用する。すると表は次のようになる。
\begin{align}
\int d\theta\,d\bar\theta\,1&=0,\\
\int d\theta\,d\bar\theta\,\theta&=0,\\
\int d\theta\,d\bar\theta\,\bar\theta&=0,\\
\int d\theta\,d\bar\theta\,\bar\theta\theta&=1.
\end{align}
したがって、
\begin{equation}
\int d\theta\,d\bar\theta\,
f(\theta,\bar\theta)=b.
\end{equation}

\begin{classroomqa}
\audiencequestion{\(\theta\)に線形な項に対して、二重積分を消えさせるのは本当に\(\theta\)の積分なのか。\lecturetimestamp{01:34:59}}
\lecturerresponse{いいえ、実際の説明では役割が逆であった。\(\theta\)を\(\theta\)について積分すると1になる。その結果には\(\bar\theta\)が含まれず、残りの1の\(\bar\theta\)積分が消える。いずれにせよ、二重積分は両方の変数が存在するときのみ生き残る。}
\end{classroomqa}

\begin{classroomqa}
\audiencequestion{\(\theta\)と\(\bar\theta\)を入れ替えると、積分は符号を変えるか。\lecturetimestamp{01:36:33}}
\lecturerresponse{はい。奇変数と奇測度は反可換である。上の規約では、
\[
\int d\bar\theta\,d\theta\,\bar\theta\theta=-1.
\]
したがって、順序を選び、それを一貫して保たなければならない。物理は規約に依存しないが、中間の符号は依存する。}
\end{classroomqa}

\(n\)個の独立なGrassmann変数に対して、パターンは直ちに得られる。
\begin{equation}
\int d\theta_n\cdots d\theta_1\,
f(\theta_1,\ldots,\theta_n)
\end{equation}
これは、変数と測度の選ばれた順序によって固定される符号を除いて、順序付けられた最高次単項式
\(\theta_1\cdots\theta_n\)の係数を取り出す。この単純な成分抽出規則が超空間作用の背後にある原動力である。超対称性の変分は奇の全微分になり、超空間積分はそれを捨てる。ちょうど通常の作用が適切な境界条件の下で時空の全微分を捨てるのと同じである。

\section{なぜこれほど小さな微積分がこれほど強力なのか}

規則はほとんど自明に見える。
\begin{equation}
\int d\theta\,1=0,
\qquad
\int d\theta\,\theta=1.
\end{equation}
しかし、これらによって、ボソンとフェルミオンの質量および結合の間の関係が自動的に保たれる相互作用場の理論を構築できる。超空間微積分の目的は、相対性理論におけるテンソル微積分の目的と同じである。すなわち、望ましい対称性が許されるあらゆる操作を通じて生き残るように記法を整えることである。

符号の帳尻合わせがその代償である。二つの奇な対象を交換するたびにマイナス符号が一つ寄与し、計算は概念的には正しくても、順序の選び方一つで失敗しうる。この困難は装飾ではない。フェルミ統計の代数的な影なのである。

\begin{classroomqa}
\audiencequestion{グラスマン数は超対称性とともに発明されたのか、それともそれ以前から存在していたのか。\lecturetimestamp{01:41:19}}
\lecturerresponse{それらはずっと古い。ヘルマン・グラスマンは19世紀にその基礎となる外積代数を展開した。微分形式はすでに反可換な基底要素を用いている。20世紀の量子場の理論と超対称性は、同じ数学に新たな物理的役割を見出した。}
\end{classroomqa}

これで具体的な構成の準備が整った。超場は完全なボソン・フェルミオン多重項をまとめ、超空間微分は代数を保ち、ベレジン積分は成分作用を抽出する。次の課題は、この仕組みを最も単純な一次元模型で用い、その後で場の理論へと持ち込むことである。
